\section{Conclusion}\label{sec:conclusion}
Sparse network--simplex convexification depends on which state information
survives the observations. On an arbitrary equality-flow graph, the cycle
rank of each unobserved subgraph gives the exact residual-coordinate count
of the observation-sensitive construction. Forest complements remove all
such coordinates. At bounded block rank, finite normal libraries also
permit original-coordinate separation and constructive recovery, including
degenerate state domains.

The coefficient results give complementary boundaries. Rank three already
requires coefficient two on a five-product $K_4$ example. A flat chain,
despite its unbounded cycle rank, admits unit flow/product coefficients
whenever at most three labels are observed, and four observed labels can
force a ratio of two. With growing label count, Fibonacci ratios occur on
simple series--parallel graphs of maximum degree three. With unrestricted
graphs, even two labels permit unbounded ratios.
Thus neither small treewidth alone nor a small number of labels alone
implies a uniform coefficient bound for the full class. These obstructions
concern original-coordinate descriptions and do not imply difficult
separation or large extension complexity.

For implementation, global state merging and native flow bounds are essential
baselines. They often make a conventional LP the fastest option. Local cycle
and observation reductions can yield further substantial size savings, while
specialized rational oracles provide exact cuts and compact decompositions
that a numerical feasibility status does not supply. The retained code,
raw measurements, and independent checks make both the advantages and the
negative runtime findings reproducible. The complete results apply to the
stated component hull; adding external coupling constraints gives a valid
relaxation and requires its own convexification analysis if an exact hull
of the enlarged model is sought.
